\section{Adaptation to unknown approximation budgets}
\label{app:approximation-adaptation}

This appendix keeps the represented-pair model of
Appendix~\ref{app:model} and strengthens the order of its estimator
quantifiers. Here the stochastic certificates $s,t$ remain known;
Appendix~\ref{app:stochastic-adaptation} subsequently removes those inputs.
No additional
geometry of either public class is assumed, and the residual variance is
bounded below only on average. The construction uses double calibration,
with a scalar correction direction as in the methodological approach of
\citet{liu2023doublecalibration}. Their sparse-GLM theorem is not invoked
for the present classes. Here a single represented propensity candidate
is restricted to a known empirical neighborhood of an initial fit; a
free vector is then chosen by constrained minimum distance. No convexity
of the public classes is needed.

\begin{theorem}[Adaptation to both approximation budgets]
\label{thm:approximation-adaptation}
Fix $(\Gamma,\rho,n,s,t)$ as in
Definition~\ref{def:indexed-compatible-laws}. There is one Borel rule
$\delta_{\Gamma,\rho,n,s,t}:\mathcal Z^{2n}\to[-3,3]$ such that,
simultaneously for every $(a,b)\in[0,1]^2$ and
$P\in\mathcal P_{\rho,n,a,s,b,t}(\Gamma)$,
\begin{equation}
 \begin{split}
 \E_P|\delta_{\Gamma,\rho,n,s,t}-\beta(P)|
 &\le C_\rho\left[q+
 \min\{a(b+t)+a^2+s^2,\ b(a+s)+(b+t)^2\}\right]\\
 &\le C_\rho(q+W).
 \end{split}
 \label{eq:approximation-adaptation-rate}
\end{equation}
The constant is independent of $\Gamma,n,s,t,a,b,P$. The rule does not
use $a,b$, a population comparator, or the covariate law. It uses the
known $s,t$. Theorem~\ref{thm:full-budget-adaptation} supplies
the separate extension that omits all four certificates.
\end{theorem}

Throughout the proof, $P_1,P_2$ denote the empirical averages on the two
original halves $\cD_1,\cD_2$, each of size $n$. Write
$f_X=(f(X_i):i\in\cD_1)$ and
$\|w\|_n^2=n^{-1}\sum_{i\in\cD_1}w_i^2$. A product inside $P_1$
can involve such vectors. Put
\begin{equation}
 q=n^{-1/2},\qquad h^2=\frac{\log(en)}n,
 \qquad\tau=n^{-12},\qquad\zeta=n^{-6}.
 \label{eq:approximation-scales}
\end{equation}
In particular $h^2\le2q$, and $\zeta,\sqrt\tau\le q$.
Every unspecified constant below depends only on $\rho$.

\subsection{Initial fits and an independent proof of the scalar pilot bound}
\label{app:approximation-pilot}

The initial treatment fit on a sample is the first represented
$\tau$-approximate minimizer of the empirical squared loss for $T$.
The initial outcome fit is the represented component of the first
$\tau$-approximate minimizer of the squared loss for $Y-\gamma T-g$,
over represented $g$ and $\gamma\in\mathbb Q\cap[-3,3]$.
These objectives use neither approximation budget. The fits on the
whole of $\cD_2$ are denoted by $\widehat\pi,\widehat\mu$.

The initial-fit bounds (B.7) and (B.9) in
Appendix~\ref{app:upper-initial}, used at $z=4\log(en)$, imply
\begin{equation}
 \begin{split}
 \E\|\widehat\pi-\pi_0\|_2^2
 &\le C_\rho(b^2+t^2+h^2),\\
 \E\|\widehat\mu-\mu_0\|_2^2
 &\le C_\rho(a^2+s^2+h^2).
 \end{split}
 \label{eq:approximation-fit-moments}
\end{equation}
Indeed their squared bounds at this one confidence level have these
orders, their failures are $O_\rho(n^{-4})$, and the fitted functions
and the true functions have fixed envelopes. The contribution of the
complements is therefore $O_\rho(n^{-4})\le C_\rho h^2$. No
unproved integration of a fitted tail is used here.

For the scalar pilot alone, split $\cD_2$ deterministically into an
internal training block of size $k=\lfloor n/2\rfloor$ and a validation
block of size $m=n-k$. Fit $g^{\rm tr},h^{\rm tr}$ on the internal
training block by the preceding rules. Let $P_{\rm va}$ be the
internal validation average, and define
\begin{align*}
 \widehat D_{\rm va}&=P_{\rm va}T(T-h^{\rm tr}),&
 \widehat A_{\rm va}&=P_{\rm va}(T-h^{\rm tr})(Y-g^{\rm tr}),\\
 \widehat V_{\rm va}&=P_{\rm va}T^2,&
 \widehat B_{\rm va}&=P_{\rm va}T(Y-g^{\rm tr}).
\end{align*}
With $\Pi=\Pi_{[-3,3]}$, the pilot is
\begin{equation}
 \bar\beta=
 \begin{cases}
  \Pi(\widehat A_{\rm va}/\widehat D_{\rm va}),
       &\widehat D_{\rm va}\ge v_-/2,\\
  \Pi\{\widehat B_{\rm va}/(\widehat V_{\rm va}\vee(v_-/2))\},
       &\widehat D_{\rm va}<v_-/2.
 \end{cases}
 \label{eq:approximation-pilot-rule}
\end{equation}
It is $\cD_2$-measurable. The full-half fits may reuse these data.

\begin{lemma}[Two second-moment bounds for one pilot]
\label{lem:approx-adapt-pilot}
For $n$ larger than a constant depending only on $\rho$, the pilot
\eqref{eq:approximation-pilot-rule} satisfies, with
$c=\beta_0-\bar\beta$,
\begin{equation}
 |c|\le6,\qquad
 \E c^2\le C_\rho
 \min\{a^2+s^2+h^2,\ b^2+t^2+h^2\}.
 \label{eq:approximation-pilot-moment}
\end{equation}
\end{lemma}
\begin{proof}
First take two fixed bounded anchors $g,h_0$ and a fresh sample of
size $m$. Let $a_0=\|\mu_0-g\|_2$, $b_0=\|\pi_0-h_0\|_2$.
Write $D,A,V,B$ for the population counterparts of the four empirical
moments in \eqref{eq:approximation-pilot-rule}. Conditional centering
gives
\[
 \begin{aligned}
 D&=v+\E\{\pi_0(\pi_0-h_0)\},&
 A&=\beta_0D+\E\{(\pi_0-h_0)(\mu_0-g)\},\\
 V&\ge v_-,& B&=\beta_0V+\E\{T(\mu_0-g)\}.
 \end{aligned}
\]
On the primary branch the clipped error is at most
\[
 \frac2{v_-}\{ |\widehat A-A|+3|\widehat D-D|+a_0b_0\}.
\]
On the fallback it is at most
\[
 \frac2{v_-}\{ |\widehat B-B|+3|\widehat V-V|+B_0a_0\},
\]
because $|V-(\widehat V\vee(v_-/2))|\le|V-\widehat V|$.
Every centered empirical moment has second moment at most $C_\rho/m$.
If $D<3v_-/4$, then $B_0b_0\ge v_-/4$, so the squared fallback
bias is bounded by $C_\rho a_0^2b_0^2$. If $D\ge3v_-/4$,
Chebyshev gives
$\Pr(\widehat D<v_-/2)\le C_\rho/m$; clipped squared loss on this
event is at most 36. Squaring the pointwise branch inequalities thus
proves
\begin{equation}
 \E|\delta_{g,h_0}-\beta_0|^2
 \le C_\rho\{m^{-1}+a_0^2b_0^2\}.
 \label{eq:approximation-anchor-mse}
\end{equation}
No independence of the branch selection and empirical moments is
required.

Condition now on the internal training block and insert its anchors.
Both anchor errors are bounded by $B_0+B_{\mathcal G}$, whence
$a_0^2b_0^2\le (B_0+B_{\mathcal G})^2\min(a_0^2,b_0^2)$.
The initial-fit moment bounds at size $k$, proved at the single
confidence level $4\log(ek)$ exactly as above, complete the proof of
\eqref{eq:approximation-pilot-moment} after averaging
\eqref{eq:approximation-anchor-mse}.

For completeness, these size-$k$ fits require no new complexity
certificate. For a countable localized difference core, adding
independent signed observations and using conditional Jensen gives
\[
 \mathfrak R_{k,P_X}(r;\Delta\cG^{\rm rep})
 \le\frac nk\mathfrak R_{n,P_X}(r;\Delta\cG^{\rm rep}).
\]
For $n\ge4$, $n/k\le3$, so the original core inequalities imply the
size-$k$ inequalities with a changed fixed process constant.
The joint-fit localization in Appendix~\ref{app:upper-initial} then
applies with the same changed constant. Moreover
$\log(ek)/k\le C h^2$ and $m^{-1}\le C h^2$.
The tolerance $\tau=n^{-12}\le k^{-12}$ is no larger than the
one permitted there. The core supremum is countable throughout; no
equality between full outer and represented localized complexities is
being asserted.
\end{proof}

\subsection{Comparison costs and the order of fixed constants}
\label{app:approximation-constants}

For a fixed compatible law choose a represented propensity comparator
$h_0$ with $\|\pi_0-h_0\|_2\le b+\zeta$. Its existence follows from
Lemma~\ref{lem:sequential-density-transfer}; it is only a proof device.
Define
\[
 B_{\rm cmp}=P_1(h_0-\widehat\pi)^2.
\]
There is a fixed $C_{\rm cmp}$, depending only on $\rho$, such that
\begin{equation}
 \Pr\{B_{\rm cmp}>C_{\rm cmp}((b+t)^2+h^2)\}
 \le C_\rho n^{-4},\qquad
 \E B_{\rm cmp}\le C_\rho((b+t)^2+h^2).
 \label{eq:approximation-comparator-cost}
\end{equation}
Here is a direct transfer proof. The $\cD_2$ treatment-fit event at
$4\log(en)$ has
$\|\widehat\pi-\pi_0\|_2^2\le C_\rho(b^2+t^2+h^2)$.
Given $\cD_2$, the function $f=\widehat\pi-h_0$ is fixed with
envelope $2B_{\mathcal G}$. Conditional scalar Bernstein,
Lemma~\ref{lem:bernstein}, gives
$P_1f^2\le2\|f\|_2^2+C_\rho h^2$, except with conditional
probability $2(en)^{-4}$. Combine this with
$\|f\|_2^2\le2\|\widehat\pi-\pi_0\|_2^2+2(b+\zeta)^2$.
This proves the tail assertion because $\zeta\le h$.
For its expectation assertion use
$\E P_1f^2=\E\|f\|_2^2$ and
\eqref{eq:approximation-fit-moments}. The selected $\cD_2$ index is
fixed under this conditioning, so no fitted main-half index is inserted
into a pointwise empirical-process event.

All constants in the procedure can now be frozen in a noncircular order.
First fix the initial-fit constants, $C_{\rm cmp}$, and the constants
of Lemma~\ref{lem:localized-process} for the two difference cores and
the bounded multipliers $u,\varepsilon,\mu_0,\pi_0$ used below.
These constants depend only on their fixed envelopes and $\rho$.
Take
\[
 R=2(B_0+B_{\mathcal G}+1),\qquad
 M=8B_0+B_{\mathcal G}+1,\qquad L=8B_0.
\]
Choose a fixed $K$ large enough that the uniform true-benchmark
process bounds in \eqref{eq:approximation-eta-event} and
\eqref{eq:approximation-true-residual-event} below are at most half
their displayed tolerances. Their proofs use only the process constants
already fixed, the bound $|c|\le6$, and quadratic absorption.
Choose $K_A\ge12C_{\rm cmp}$ and put
\begin{equation}
 \kappa_\mu=K(s^2+h^2),\qquad
 \kappa_\pi=K(t^2+h^2),\qquad A_t=K_A(t^2+h^2).
 \label{eq:approximation-tolerances}
\end{equation}
Finally choose $n_\rho\ge4$ so that all the confidence levels
$4\log(en)$, and their size-$k$ counterparts, are in the process
tail ranges, and the scalar denominator event below has its stated
margin for $n\ge n_\rho$. This last choice changes none of the other
constants. All these choices precede the budgets, pair, law, and data.
For $n<n_\rho$ define the rule to be zero, whose loss is absorbed by
$C_\rho q$. In what follows $n\ge n_\rho$.

\subsection{The two calibration programs}
\label{app:approximation-programs}

Let
\[
 \mathcal F_X=\sigma(\cD_2,X_i:i\in\cD_1),\qquad
 \mathcal F_T=\sigma(\mathcal F_X,T_i:i\in\cD_1).
\]
The first program constructs an outcome correction without using the
main-half treatments or outcomes. For an $n$-vector $m$ define
\[
 H^{\rm pil}_\pi(m)=
 \sup_{f\in\Delta\cG_\pi^{\rm rep}}
 \{|P_2[(Y-\bar\beta T)f]-P_1(mf)|-P_1f^2\}.
\]
Retain $m\in\mathbb Q^n$ satisfying
\begin{equation}
 H^{\rm pil}_\pi(m)<\kappa_\pi,\qquad \|m\|_n<M.
 \label{eq:approximation-eta-program}
\end{equation}
If this set is nonempty, let $\widehat\eta$ be its first point whose
$\|m-\widehat\mu_X\|_n^2$ is less than the feasible countable
infimum plus $\tau$. Otherwise output zero and set $\widehat\eta=0$
for analytical vector conventions. Write
$d=\widehat\eta-\widehat\mu_X$. Both vectors are
$\mathcal F_X$-measurable.

For the second program write $(h_j)_{j\ge1}$ for the supplied
propensity core. Enumerate all pairs $(j,z)\in\mathbb N\times\mathbb Q^n$.
For each pair set
\[
 \pi^c=(h_j)_X+z,\qquad r=T_X-\pi^c,
 \qquad D=P_1(Tr),\qquad B_j=P_1(h_j-\widehat\pi)^2.
\]
For $\nu\in\{\mu,\pi\}$ define
$H_\nu(r)=\sup_{f\in\Delta\cG_\nu^{\rm rep}}
\{|P_1(rf)|-P_1f^2\}$. Retain exactly the pairs satisfying
\begin{equation}
 \begin{gathered}
 B_j<A_t,\qquad D>v_-/2,\qquad
 \|r\|_n<R,\qquad\|\pi^c\|_n<R,\\
 H_\mu(r)<\kappa_\mu,\qquad H_\pi(r)<\kappa_\pi,\\
 |P_1(rd)|<Lh\|d\|_n+\tau,\\
 |P_1(r\pi^c)|<Lh\|\pi^c\|_n+\tau,\\
 |P_1(rz)|<Lh\|z\|_n+\tau.
 \end{gathered}
 \label{eq:approximation-propensity-program}
\end{equation}
Only the free-vector objective $\|z\|_n^2$ is minimized: choose the
first retained pair whose objective is less than its countable feasible
infimum plus $\tau$. Denote the selected candidate function by
$\widehat h$, and its other quantities by
$\widehat z,\widehat\pi^c,\widehat r,\widehat D$.
If either program is empty, output zero; define $\widehat r=0$ there
for the conditional noise arguments. Otherwise output
\begin{equation}
 \delta=\Pi_{[-3,3]}\left[
 \bar\beta+
 \frac{P_1\{\widehat r(Y-\bar\beta T-\widehat\eta)\}}
      {\widehat D}\right]
 =\Pi_{[-3,3]}
 \frac{P_1\{\widehat r(Y-\widehat\eta)\}}{\widehat D}.
 \label{eq:approximation-rule}
\end{equation}
The denominator is positive on every feasible sample. The second
equality follows from its definition, not from an approximate moment
identity. Neither program uses $a,b$.

Each program is Borel. The cores are countable; all suprema in the
constraints are countable and are continuous in the calibrated vector
by the common difference envelope. All other constraints are continuous
in the finite-dimensional vector, for a fixed represented index.
Countable feasible-set decisions, countable infima, and first strict
approximate minimizers are consequently Borel. The strict inequalities
and the scalar $+\tau$ terms provide open neighborhoods even when a
calibrated direction is zero. Thus a strictly feasible real vector can
be approximated by retained rational vectors with objective converging
to its own objective. No exact minimizer is required. The initial fits
and the scalar pilot are Borel by their countable selection conventions
and explicit ratios, respectively. In particular $\widehat r$ is
$\mathcal F_T$-measurable: no main-half outcome is used to choose it.

\subsection{Feasibility, margins, and two comparisons for the same program}
\label{app:approximation-feasibility}

Fix a law and represented comparators $g,h_0$ whose approximation
errors are at most $a+\zeta,b+\zeta$, respectively. Write
\[
 \eta_0=(\mu_0+c\pi_0)_X,\quad
 r_\eta=\widehat\eta-\eta_0,\quad
 e_\pi=(\widehat\pi-\pi_0)_X,\quad
 r_\pi=(\pi_0-h_0)_X.
\]
The random function defining $\eta_0$ depends on $\cD_2$ only
through the bounded scalar $c$; its vector also uses the main covariates.
Before inserting that scalar,
form uniform process events using the exact identity
\begin{equation}
 \begin{split}
 P_2[(Y-\bar\beta T)f]-P_1(\eta_0f)
 ={}&P_2(\varepsilon f)+cP_2(uf)\\
 &+(P_2-P_1)(\mu_0f)+c(P_2-P_1)(\pi_0f).
 \end{split}
 \label{eq:approximation-eta-identity}
\end{equation}
Lemma~\ref{lem:localized-process}, applied to the fixed bounded
multipliers displayed here and to the square process, gives a finite
intersection of events, of failure $O_\rho(n^{-4})$, on which the
absolute value of the right side minus $P_1f^2$ is bounded by
\[
 C_\rho(t+h)\|f\|_2-\|f\|_2^2+C_\rho(t^2+h^2)
 \le C_\rho(t^2+h^2)
\]
uniformly over the propensity difference core. The chosen $K$ therefore
supplies the event
\begin{equation}
 G_\eta=\{H^{\rm pil}_\pi(\eta_0)\le\kappa_\pi/2\}
 \in\mathcal F_X,\qquad
 \Pr(G_\eta^c)\le C_\rho n^{-4}.
 \label{eq:approximation-eta-event}
\end{equation}
Also $\|\eta_0\|_n\le7B_0<M$, with strict norm margin.
On $G_\eta$ the first program is consequently nonempty and its
comparison with $\eta_0$ yields
\begin{equation}
 \begin{gathered}
 \|d\|_n\le\|\eta_0-\widehat\mu_X\|_n+\sqrt\tau,\qquad
 \|r_\eta\|_n\le2\|\eta_0-\widehat\mu_X\|_n+\sqrt\tau,\\
 |P_1(r_\eta f)|\le2P_1f^2+3\kappa_\pi/2
 \quad(f\in\Delta\cG_\pi^{\rm rep}).
 \end{gathered}
 \label{eq:approximation-eta-comparison}
\end{equation}
The vector norms of $d,r_\eta$ are bounded by fixed constants even
outside this event, using the norm constraint and the fallback
convention. Moreover \eqref{eq:approximation-fit-moments}, the first
pilot moment bound, and independence of the main covariates give
\begin{equation}
 \E1_{G_\eta}(\|d\|_n^2+\|r_\eta\|_n^2)
 \le C_\rho(a^2+s^2+h^2).
 \label{eq:approximation-eta-moments}
\end{equation}
Independence between $\bar\beta$ and the full-half fits is unnecessary.
The same moment bound also holds without $1_{G_\eta}$: both vector
norms are uniformly bounded on its complement, whose probability is
$O_\rho(n^{-4})\le C_\rho h^2$.

The same process lemma and quadratic absorption, now with multiplier
$u$ on the main half, give
\begin{equation}
 H_\mu(u_X)\le\kappa_\mu/2,\qquad
 H_\pi(u_X)\le\kappa_\pi/2
 \label{eq:approximation-true-residual-event}
\end{equation}
except with probability $C_\rho n^{-4}$. For scalar constraints,
conditionally on $\mathcal F_X$ and any of the four vectors
\[
 w=d,\quad w=(\pi_0)_X,\quad
 w=-e_\pi=(\pi_0-\widehat\pi)_X,\quad w=r_\pi,
\]
weighted Hoeffding gives, at confidence $4\log(en)$,
\begin{equation}
 |P_1(uw)|\le\sqrt8 B_0h\|w\|_n
 \le (L/2)h\|w\|_n,
 \label{eq:approximation-scalar-margin}
\end{equation}
with conditional failure at most $2(en)^{-4}$ for each vector.
These statements also hold for the random vector $-e_\pi$: conditioning
fixes the full-half fit and main covariates, and leaves the $u_i$
independent, centered, and bounded. No coordinate bound on $w$ is
needed for weighted Hoeffding. If $w=0$, the added $\tau>0$ in the
program still makes its scalar constraint strict.

Finally bounded scalar Bernstein and $\E(Tu)=v\ge v_-$ yield
$P_1(Tu)>3v_-/4$ except with probability $C_\rho n^{-4}$, by the
choice of $n_\rho$. Let $G$ be the finite intersection of $G_\eta$,
the two events in \eqref{eq:approximation-true-residual-event}, the four
scalar events \eqref{eq:approximation-scalar-margin}, this denominator
event, and the comparison-cost event in
\eqref{eq:approximation-comparator-cost}. Then
\begin{equation}
 G\in\mathcal F_T,\qquad \Pr(G^c)\le C_\rho n^{-4}.
 \label{eq:approximation-good-event}
\end{equation}
All uniform class events precede every fitted-index substitution.

On $G$, consider the following two representations of the same real
vector $(\pi_0)_X$ in the second program. First take the represented
candidate $h_j=\widehat\pi$ and $z=-e_\pi$. Its comparison cost is
zero, and $\pi^c=(\pi_0)_X$, $r=u_X$. All its constraints have the
strict margins just proved, including
$\|u\|_n,\|\pi_0\|_n\le B_0<R$. Thus the second program is
nonempty for every $a,b$.

For the other comparison suppose only in the proof that $b\le t+h$.
Then
\[
 (b+t)^2+h^2\le(2t+h)^2+h^2\le6(t^2+h^2).
\]
The cost event and $K_A\ge12C_{\rm cmp}$ give
\begin{equation}
 B_{\rm cmp}\le A_t/2.
 \label{eq:approximation-local-comparator}
\end{equation}
Take the fixed candidate $h_j=h_0$ and $z=r_\pi$. Again
$\pi^c=(\pi_0)_X$, $r=u_X$, and every constraint is strictly
feasible. In each comparison the represented index is a legitimate
index of the program, and its real free vector can be approximated
by rational vectors. Minimization of the same free-vector objective
therefore gives, on $G$,
\begin{equation}
 \begin{aligned}
 \|\widehat z\|_n&\le\|e_\pi\|_n+\sqrt\tau
       &&\text{for all }b,\\
 \|\widehat z\|_n&\le\|r_\pi\|_n+\sqrt\tau
       &&\text{if }b\le t+h.
 \end{aligned}
 \label{eq:approximation-two-comparisons}
\end{equation}
The algorithm does not test this condition and does not choose between
the two analytical comparisons.

Put $\Delta=\widehat\pi^c-(\pi_0)_X$. The empirical neighborhood
constraint and the first line of
\eqref{eq:approximation-two-comparisons} imply
\[
 \|\Delta\|_n
 \le\|e_\pi\|_n+\|(\widehat h-\widehat\pi)_X\|_n
       +\|\widehat z\|_n
 \le2\|e_\pi\|_n+\sqrt{A_t}+\sqrt\tau.
\]
Consequently
\begin{equation}
 \E1_G\|\Delta\|_n^2\le C_\rho(b^2+t^2+h^2).
 \label{eq:approximation-propensity-moment}
\end{equation}
On every feasible sample,
$\|\widehat z\|_n\le\|\widehat\pi^c\|_n+\|\widehat h_X\|_n
<R+B_{\mathcal G}$; this uniform bound is useful for scalar terms.

\subsection{Two outcome analyses of the same rule}
\label{app:approximation-outcome-bound}

Conditional on $\mathcal F_T$, the $\varepsilon_i$ are independent
and centered, and $\|\widehat r\|_n<R$ on feasible samples. Thus
\begin{equation}
 \E|P_1(\varepsilon\widehat r)|\le C_\rho q.
 \label{eq:approximation-outcome-noise}
\end{equation}
Likewise conditioning on $\mathcal F_X$ gives $C_\rho q$ bounds for
$P_1(ur_\eta)$, $P_1[u(\mu_0-\widehat\mu)]$, and
$P_1[u(\mu_0-g)]$. Their empirical coefficient norms are uniformly
bounded. Good-event indicators can first be dropped from nonnegative
absolute values. The denominator can be bounded below pointwise by
$v_-/2$ before applying these bounds. No conditioning on a selected
main-half outcome is used.

On $G$ the raw, unclipped output has the exact identity
\begin{equation}
 \widehat D(\delta^{\rm raw}-\beta_0)
 =P_1(\varepsilon\widehat r)
  +P_1\{(\mu_0-\widehat\mu)\widehat r\}
  -P_1(d\widehat r).
 \label{eq:approximation-outcome-identity}
\end{equation}
Split $\mu_0-\widehat\mu=(\mu_0-g)+(g-\widehat\mu)$ and
use $\widehat r=u_X-\Delta$. The centered fixed-comparator term
costs $C_\rho q$. Cauchy--Schwarz and
\eqref{eq:approximation-propensity-moment} bound the other
fixed-comparator term by $C_\rho(a+\zeta)(b+t+h)$ in expectation.
The outcome-class constraint in
\eqref{eq:approximation-propensity-program}, uniform before the
selected represented index is inserted, gives
\[
 |P_1\{(g-\widehat\mu)\widehat r\}|
 \le P_1(g-\widehat\mu)^2+\kappa_\mu.
\]
Its expectation is at most $C_\rho(a^2+s^2+h^2)$ by the initial
fit moment bound and the fixed-comparator error. The correction
constraint and \eqref{eq:approximation-eta-moments} give
$\E1_G|P_1(d\widehat r)|\le C_\rho h(a+s+h)+\tau$.
Use $ah\le(a^2+h^2)/2$, $sh\le(s^2+h^2)/2$, and $h^2\le2q$.
Clipping cannot increase the error, and clipped loss on $G^c$ is at
most six. We obtain
\begin{equation}
 \E|\delta-\beta_0|
 \le C_\rho\{q+a(b+t)+a^2+s^2\}.
 \label{eq:approximation-outcome-rate}
\end{equation}

Keeping $\mu_0-\widehat\mu$ together instead in
\eqref{eq:approximation-outcome-identity} gives a second bound for
this very same rule. The centered term again costs $C_\rho q$;
Cauchy--Schwarz bounds the $\Delta$ term by
$C_\rho(a+s+h)(b+t+h)$, and the correction term by
$C_\rho h(a+s+h)+\tau$. Therefore
\begin{equation}
 \E|\delta-\beta_0|
 \le C_\rho\{q+(a+s+h)(b+t+2h)\}.
 \label{eq:approximation-product-rate}
\end{equation}
The floor $h$ has not been replaced by $q$ in this display.

\subsection{The propensity analysis and the proof of the theorem}
\label{app:approximation-propensity-bound}

If $b>t+h$, the product bound already gives the required propensity
rate: $b+t+2h\le4b$ and
$bh\le(b^2+h^2)/2$. Thus
\begin{equation}
 \E|\delta-\beta_0|
 \le C_\rho\{q+b(a+s)+(b+t)^2\}
 \qquad(b>t+h).
 \label{eq:approximation-propensity-large-b}
\end{equation}
It remains to analyze $b\le t+h$, where the second comparison in
\eqref{eq:approximation-two-comparisons} is available.

Because $\mu_0-\widehat\eta=-r_\eta-c\pi_0$, direct substitution
in \eqref{eq:approximation-rule} gives
\begin{equation}
 \widehat D(\delta^{\rm raw}-\beta_0)
 =P_1(\varepsilon\widehat r)-P_1(r_\eta\widehat r)
       -cP_1(\pi_0\widehat r).
 \label{eq:approximation-propensity-identity}
\end{equation}
Let $f_1=(\widehat h-\widehat\pi)_X$ and
$f_0=(\widehat\pi-h_0)_X$. Both are represented propensity
differences, and
$\Delta=f_1+\widehat z+f_0-r_\pi$. The two calibration
inequalities in \eqref{eq:approximation-eta-comparison} give
\begin{equation}
 \begin{split}
 |P_1(r_\eta\Delta)|
 \le{}&2A_t+2B_{\rm cmp}+3\kappa_\pi\\
 &+\|r_\eta\|_n(\|\widehat z\|_n+\|r_\pi\|_n).
 \end{split}
 \label{eq:approximation-remainder-calibration}
\end{equation}
The second comparison for $\widehat z$,
\eqref{eq:approximation-eta-moments}, and
$\E\|r_\pi\|_n^2\le(b+\zeta)^2$ bound its good-event
expectation by
\begin{equation}
 C_\rho\{(b+t)^2+h^2+(b+\zeta)(a+s+h)+\sqrt\tau\}.
 \label{eq:approximation-remainder-bound}
\end{equation}
The term $P_1(r_\eta\widehat r)
=P_1(r_\eta u)-P_1(r_\eta\Delta)$ has only an additional
$C_\rho q$ contribution, by the conditional noise bound.

For the final term in
\eqref{eq:approximation-propensity-identity}, use
$\pi_0=\widehat\pi^c-\Delta$. The propensity difference
constraints and the two scalar constraints for
$\widehat\pi^c,\widehat z$ give, pointwise on $G$,
\begin{equation}
 \begin{split}
 |P_1(\pi_0\widehat r)|
 \le{}&Lh\|\widehat\pi^c\|_n+\tau
       +A_t+B_{\rm cmp}+2\kappa_\pi\\
 &+Lh\|\widehat z\|_n+\tau+R\|r_\pi\|_n.
 \end{split}
 \label{eq:approximation-self-calibration}
\end{equation}
Multiply by $|c|$. Bounded $|c|$, the deterministic cap, and
\eqref{eq:approximation-comparator-cost} control the cap, comparator,
and tolerance terms by $C_\rho((b+t)^2+h^2)$ in expectation.
Uniform norm bounds reduce the two terms proportional to $h$ to
$C_\rho h\E|c|$. The second bound on the one pilot in
\eqref{eq:approximation-pilot-moment} gives
\begin{equation}
 h\E|c|\le C_\rho h(b+t+h)
 \le C_\rho((b+t)^2+h^2).
 \label{eq:approximation-pilot-use-propensity}
\end{equation}
Its first bound, together with Cauchy--Schwarz, gives
\begin{equation}
 \E|c|\|r_\pi\|_n
 \le C_\rho(b+\zeta)(a+s+h).
 \label{eq:approximation-pilot-use-outcome}
\end{equation}
The pilot and fits may be dependent; these bounds do not condition
their shared fitting data on a fitted coefficient.

Combine \eqref{eq:approximation-propensity-identity}--
\eqref{eq:approximation-pilot-use-outcome}, divide by the imposed
denominator bound, and use $bh\le(b^2+h^2)/2$, $h^2\le2q$,
and $\zeta,\sqrt\tau\le q$. Clipping and
\eqref{eq:approximation-good-event} add at most
$6\Pr(G^c)=O_\rho(n^{-4})$. It follows that
\begin{equation}
 \E|\delta-\beta_0|
 \le C_\rho\{q+b(a+s)+(b+t)^2\}
 \qquad(b\le t+h).
 \label{eq:approximation-propensity-small-b}
\end{equation}
The two cases \eqref{eq:approximation-propensity-large-b} and
\eqref{eq:approximation-propensity-small-b} partition only the proof.
Their programs, constants, and outputs are identical. Together with
\eqref{eq:approximation-outcome-rate}, they prove
\eqref{eq:approximation-adaptation-rate} for all $n\ge n_\rho$.
The zero output covers smaller $n$, since its loss is at most three
and $q\ge n_\rho^{-1/2}$. Positive $h$ makes the tolerances and
empirical neighborhood well defined when $s=0$ or $t=0$. Fixed
$\zeta$-comparators exist even if an approximation infimum is not
attained, including when $a=0$ or $b=0$; their remaining terms were
explicitly absorbed above. This proves
Theorem~\ref{thm:approximation-adaptation}. \hfill$\square$
